\documentclass[11pt,a4paper]{article}
% Report v3 — every number comes from results/macros.tex (written by ser_tolerance_v3).
% Missing numbers render as a bold [key] so nothing unmeasured can slip into the PDF.
\usepackage[margin=2.2cm]{geometry}
\usepackage{amsmath,amssymb,graphicx,booktabs,xcolor,enumitem,float,microtype,hyperref,caption}
\hypersetup{colorlinks=true,linkcolor=blue!50!black,citecolor=blue!50!black,urlcolor=blue!50!black}
\captionsetup{font=small,labelfont=bf}
\setlist{leftmargin=*,itemsep=2pt,topsep=3pt}
\emergencystretch=2em

\makeatletter
\newcommand{\res}[1]{\ifcsname res@#1\endcsname\csname res@#1\endcsname\else\textbf{\textcolor{red!70!black}{[\texttt{\detokenize{#1}}]}}\fi}
\makeatother
\IfFileExists{results/macros.tex}{% auto-generated by ser_tolerance_v3 — do not edit by hand
\expandafter\def\csname res@n_trials\endcsname{496}
\expandafter\def\csname res@n_neurons\endcsname{161}
\expandafter\def\csname res@n_m1\endcsname{67}
\expandafter\def\csname res@n_pmd\endcsname{94}
\expandafter\def\csname res@sparsity\endcsname{93.9}
\expandafter\def\csname res@rate_hz\endcsname{6.41}
\expandafter\def\csname res@kbar\endcsname{9.8}
\expandafter\def\csname res@fhz_spike_1e-4\endcsname{0.000321}
\expandafter\def\csname res@fhz_1ms_1e-4\endcsname{0.0994}
\expandafter\def\csname res@fhz_spike_3e-4\endcsname{0.000962}
\expandafter\def\csname res@fhz_1ms_3e-4\endcsname{0.298}
\expandafter\def\csname res@fhz_spike_1e-3\endcsname{0.00321}
\expandafter\def\csname res@fhz_1ms_1e-3\endcsname{0.994}
\expandafter\def\csname res@fhz_spike_3e-3\endcsname{0.00962}
\expandafter\def\csname res@fhz_1ms_3e-3\endcsname{2.98}
\expandafter\def\csname res@fhz_spike_1e-2\endcsname{0.0321}
\expandafter\def\csname res@fhz_1ms_1e-2\endcsname{9.94}
\expandafter\def\csname res@fhz_spike_3e-2\endcsname{0.0962}
\expandafter\def\csname res@fhz_1ms_3e-2\endcsname{29.8}
\expandafter\def\csname res@fhz_spike_1e-1\endcsname{0.321}
\expandafter\def\csname res@fhz_1ms_1e-1\endcsname{99.4}
\expandafter\def\csname res@r_Ridge\endcsname{0.858}
\expandafter\def\csname res@r_Ridge_ci\endcsname{[0.852, 0.865]}
\expandafter\def\csname res@rvx_Ridge\endcsname{0.858}
\expandafter\def\csname res@rsm_Ridge\endcsname{0.869}
\expandafter\def\csname res@r_MLP\endcsname{0.907}
\expandafter\def\csname res@r_MLP_ci\endcsname{[0.902, 0.912]}
\expandafter\def\csname res@rvx_MLP\endcsname{0.909}
\expandafter\def\csname res@rsm_MLP\endcsname{0.922}
\expandafter\def\csname res@r_LSTM\endcsname{0.903}
\expandafter\def\csname res@r_LSTM_ci\endcsname{[0.898, 0.910]}
\expandafter\def\csname res@rvx_LSTM\endcsname{0.904}
\expandafter\def\csname res@rsm_LSTM\endcsname{0.909}
\expandafter\def\csname res@r_Transformer\endcsname{0.896}
\expandafter\def\csname res@r_Transformer_ci\endcsname{[0.890, 0.903]}
\expandafter\def\csname res@rvx_Transformer\endcsname{0.899}
\expandafter\def\csname res@rsm_Transformer\endcsname{0.905}
\expandafter\def\csname res@rho_Ridge_per_spike_0\endcsname{1.000}
\expandafter\def\csname res@rho_Ridge_per_spike_0_ci\endcsname{[1.000, 1.000]}
\expandafter\def\csname res@rho_Ridge_per_spike_1e-4\endcsname{1.000}
\expandafter\def\csname res@rho_Ridge_per_spike_1e-4_ci\endcsname{[1.000, 1.000]}
\expandafter\def\csname res@sig_Ridge_per_spike_1e-4\endcsname{no}
\expandafter\def\csname res@rho_Ridge_per_spike_3e-4\endcsname{1.000}
\expandafter\def\csname res@rho_Ridge_per_spike_3e-4_ci\endcsname{[1.000, 1.000]}
\expandafter\def\csname res@sig_Ridge_per_spike_3e-4\endcsname{no}
\expandafter\def\csname res@rho_Ridge_per_spike_1e-3\endcsname{1.000}
\expandafter\def\csname res@rho_Ridge_per_spike_1e-3_ci\endcsname{[1.000, 1.000]}
\expandafter\def\csname res@sig_Ridge_per_spike_1e-3\endcsname{yes}
\expandafter\def\csname res@rho_Ridge_per_spike_3e-3\endcsname{1.000}
\expandafter\def\csname res@rho_Ridge_per_spike_3e-3_ci\endcsname{[0.999, 1.000]}
\expandafter\def\csname res@sig_Ridge_per_spike_3e-3\endcsname{yes}
\expandafter\def\csname res@rho_Ridge_per_spike_1e-2\endcsname{0.998}
\expandafter\def\csname res@rho_Ridge_per_spike_1e-2_ci\endcsname{[0.998, 0.999]}
\expandafter\def\csname res@sig_Ridge_per_spike_1e-2\endcsname{yes}
\expandafter\def\csname res@rho_Ridge_per_spike_3e-2\endcsname{0.995}
\expandafter\def\csname res@rho_Ridge_per_spike_3e-2_ci\endcsname{[0.995, 0.996]}
\expandafter\def\csname res@sig_Ridge_per_spike_3e-2\endcsname{yes}
\expandafter\def\csname res@rho_Ridge_per_spike_1e-1\endcsname{0.983}
\expandafter\def\csname res@rho_Ridge_per_spike_1e-1_ci\endcsname{[0.981, 0.984]}
\expandafter\def\csname res@sig_Ridge_per_spike_1e-1\endcsname{yes}
\expandafter\def\csname res@rho_Ridge_per_1ms_0\endcsname{1.000}
\expandafter\def\csname res@rho_Ridge_per_1ms_0_ci\endcsname{[1.000, 1.000]}
\expandafter\def\csname res@rho_Ridge_per_1ms_1e-4\endcsname{0.996}
\expandafter\def\csname res@rho_Ridge_per_1ms_1e-4_ci\endcsname{[0.996, 0.997]}
\expandafter\def\csname res@sig_Ridge_per_1ms_1e-4\endcsname{yes}
\expandafter\def\csname res@rho_Ridge_per_1ms_3e-4\endcsname{0.989}
\expandafter\def\csname res@rho_Ridge_per_1ms_3e-4_ci\endcsname{[0.988, 0.990]}
\expandafter\def\csname res@sig_Ridge_per_1ms_3e-4\endcsname{yes}
\expandafter\def\csname res@rho_Ridge_per_1ms_1e-3\endcsname{0.964}
\expandafter\def\csname res@rho_Ridge_per_1ms_1e-3_ci\endcsname{[0.963, 0.966]}
\expandafter\def\csname res@sig_Ridge_per_1ms_1e-3\endcsname{yes}
\expandafter\def\csname res@rho_Ridge_per_1ms_3e-3\endcsname{0.901}
\expandafter\def\csname res@rho_Ridge_per_1ms_3e-3_ci\endcsname{[0.898, 0.905]}
\expandafter\def\csname res@sig_Ridge_per_1ms_3e-3\endcsname{yes}
\expandafter\def\csname res@rho_Ridge_per_1ms_1e-2\endcsname{0.756}
\expandafter\def\csname res@rho_Ridge_per_1ms_1e-2_ci\endcsname{[0.747, 0.764]}
\expandafter\def\csname res@sig_Ridge_per_1ms_1e-2\endcsname{yes}
\expandafter\def\csname res@rho_Ridge_per_1ms_3e-2\endcsname{0.543}
\expandafter\def\csname res@rho_Ridge_per_1ms_3e-2_ci\endcsname{[0.535, 0.551]}
\expandafter\def\csname res@sig_Ridge_per_1ms_3e-2\endcsname{yes}
\expandafter\def\csname res@rho_Ridge_per_1ms_1e-1\endcsname{0.302}
\expandafter\def\csname res@rho_Ridge_per_1ms_1e-1_ci\endcsname{[0.290, 0.316]}
\expandafter\def\csname res@sig_Ridge_per_1ms_1e-1\endcsname{yes}
\expandafter\def\csname res@rho_MLP_per_spike_0\endcsname{1.000}
\expandafter\def\csname res@rho_MLP_per_spike_0_ci\endcsname{[1.000, 1.000]}
\expandafter\def\csname res@rho_MLP_per_spike_1e-4\endcsname{1.000}
\expandafter\def\csname res@rho_MLP_per_spike_1e-4_ci\endcsname{[1.000, 1.000]}
\expandafter\def\csname res@sig_MLP_per_spike_1e-4\endcsname{no}
\expandafter\def\csname res@rho_MLP_per_spike_3e-4\endcsname{1.000}
\expandafter\def\csname res@rho_MLP_per_spike_3e-4_ci\endcsname{[1.000, 1.000]}
\expandafter\def\csname res@sig_MLP_per_spike_3e-4\endcsname{no}
\expandafter\def\csname res@rho_MLP_per_spike_1e-3\endcsname{1.000}
\expandafter\def\csname res@rho_MLP_per_spike_1e-3_ci\endcsname{[1.000, 1.000]}
\expandafter\def\csname res@sig_MLP_per_spike_1e-3\endcsname{no}
\expandafter\def\csname res@rho_MLP_per_spike_3e-3\endcsname{0.999}
\expandafter\def\csname res@rho_MLP_per_spike_3e-3_ci\endcsname{[0.999, 1.000]}
\expandafter\def\csname res@sig_MLP_per_spike_3e-3\endcsname{yes}
\expandafter\def\csname res@rho_MLP_per_spike_1e-2\endcsname{0.998}
\expandafter\def\csname res@rho_MLP_per_spike_1e-2_ci\endcsname{[0.997, 0.998]}
\expandafter\def\csname res@sig_MLP_per_spike_1e-2\endcsname{yes}
\expandafter\def\csname res@rho_MLP_per_spike_3e-2\endcsname{0.993}
\expandafter\def\csname res@rho_MLP_per_spike_3e-2_ci\endcsname{[0.992, 0.993]}
\expandafter\def\csname res@sig_MLP_per_spike_3e-2\endcsname{yes}
\expandafter\def\csname res@rho_MLP_per_spike_1e-1\endcsname{0.976}
\expandafter\def\csname res@rho_MLP_per_spike_1e-1_ci\endcsname{[0.972, 0.979]}
\expandafter\def\csname res@sig_MLP_per_spike_1e-1\endcsname{yes}
\expandafter\def\csname res@rho_MLP_per_1ms_0\endcsname{1.000}
\expandafter\def\csname res@rho_MLP_per_1ms_0_ci\endcsname{[1.000, 1.000]}
\expandafter\def\csname res@rho_MLP_per_1ms_1e-4\endcsname{0.994}
\expandafter\def\csname res@rho_MLP_per_1ms_1e-4_ci\endcsname{[0.993, 0.994]}
\expandafter\def\csname res@sig_MLP_per_1ms_1e-4\endcsname{yes}
\expandafter\def\csname res@rho_MLP_per_1ms_3e-4\endcsname{0.982}
\expandafter\def\csname res@rho_MLP_per_1ms_3e-4_ci\endcsname{[0.979, 0.984]}
\expandafter\def\csname res@sig_MLP_per_1ms_3e-4\endcsname{yes}
\expandafter\def\csname res@rho_MLP_per_1ms_1e-3\endcsname{0.942}
\expandafter\def\csname res@rho_MLP_per_1ms_1e-3_ci\endcsname{[0.936, 0.948]}
\expandafter\def\csname res@sig_MLP_per_1ms_1e-3\endcsname{yes}
\expandafter\def\csname res@rho_MLP_per_1ms_3e-3\endcsname{0.838}
\expandafter\def\csname res@rho_MLP_per_1ms_3e-3_ci\endcsname{[0.828, 0.847]}
\expandafter\def\csname res@sig_MLP_per_1ms_3e-3\endcsname{yes}
\expandafter\def\csname res@rho_MLP_per_1ms_1e-2\endcsname{0.602}
\expandafter\def\csname res@rho_MLP_per_1ms_1e-2_ci\endcsname{[0.591, 0.613]}
\expandafter\def\csname res@sig_MLP_per_1ms_1e-2\endcsname{yes}
\expandafter\def\csname res@rho_MLP_per_1ms_3e-2\endcsname{0.360}
\expandafter\def\csname res@rho_MLP_per_1ms_3e-2_ci\endcsname{[0.350, 0.370]}
\expandafter\def\csname res@sig_MLP_per_1ms_3e-2\endcsname{yes}
\expandafter\def\csname res@rho_MLP_per_1ms_1e-1\endcsname{0.176}
\expandafter\def\csname res@rho_MLP_per_1ms_1e-1_ci\endcsname{[0.168, 0.185]}
\expandafter\def\csname res@sig_MLP_per_1ms_1e-1\endcsname{yes}
\expandafter\def\csname res@rho_LSTM_per_spike_0\endcsname{1.000}
\expandafter\def\csname res@rho_LSTM_per_spike_0_ci\endcsname{[1.000, 1.000]}
\expandafter\def\csname res@rho_LSTM_per_spike_1e-4\endcsname{1.000}
\expandafter\def\csname res@rho_LSTM_per_spike_1e-4_ci\endcsname{[1.000, 1.000]}
\expandafter\def\csname res@sig_LSTM_per_spike_1e-4\endcsname{no}
\expandafter\def\csname res@rho_LSTM_per_spike_3e-4\endcsname{1.000}
\expandafter\def\csname res@rho_LSTM_per_spike_3e-4_ci\endcsname{[1.000, 1.000]}
\expandafter\def\csname res@sig_LSTM_per_spike_3e-4\endcsname{no}
\expandafter\def\csname res@rho_LSTM_per_spike_1e-3\endcsname{1.000}
\expandafter\def\csname res@rho_LSTM_per_spike_1e-3_ci\endcsname{[1.000, 1.000]}
\expandafter\def\csname res@sig_LSTM_per_spike_1e-3\endcsname{no}
\expandafter\def\csname res@rho_LSTM_per_spike_3e-3\endcsname{0.999}
\expandafter\def\csname res@rho_LSTM_per_spike_3e-3_ci\endcsname{[0.999, 1.000]}
\expandafter\def\csname res@sig_LSTM_per_spike_3e-3\endcsname{no}
\expandafter\def\csname res@rho_LSTM_per_spike_1e-2\endcsname{0.998}
\expandafter\def\csname res@rho_LSTM_per_spike_1e-2_ci\endcsname{[0.998, 0.999]}
\expandafter\def\csname res@sig_LSTM_per_spike_1e-2\endcsname{yes}
\expandafter\def\csname res@rho_LSTM_per_spike_3e-2\endcsname{0.994}
\expandafter\def\csname res@rho_LSTM_per_spike_3e-2_ci\endcsname{[0.993, 0.994]}
\expandafter\def\csname res@sig_LSTM_per_spike_3e-2\endcsname{yes}
\expandafter\def\csname res@rho_LSTM_per_spike_1e-1\endcsname{0.978}
\expandafter\def\csname res@rho_LSTM_per_spike_1e-1_ci\endcsname{[0.976, 0.980]}
\expandafter\def\csname res@sig_LSTM_per_spike_1e-1\endcsname{yes}
\expandafter\def\csname res@rho_LSTM_per_1ms_0\endcsname{1.000}
\expandafter\def\csname res@rho_LSTM_per_1ms_0_ci\endcsname{[1.000, 1.000]}
\expandafter\def\csname res@rho_LSTM_per_1ms_1e-4\endcsname{0.995}
\expandafter\def\csname res@rho_LSTM_per_1ms_1e-4_ci\endcsname{[0.994, 0.996]}
\expandafter\def\csname res@sig_LSTM_per_1ms_1e-4\endcsname{yes}
\expandafter\def\csname res@rho_LSTM_per_1ms_3e-4\endcsname{0.985}
\expandafter\def\csname res@rho_LSTM_per_1ms_3e-4_ci\endcsname{[0.984, 0.987]}
\expandafter\def\csname res@sig_LSTM_per_1ms_3e-4\endcsname{yes}
\expandafter\def\csname res@rho_LSTM_per_1ms_1e-3\endcsname{0.948}
\expandafter\def\csname res@rho_LSTM_per_1ms_1e-3_ci\endcsname{[0.943, 0.952]}
\expandafter\def\csname res@sig_LSTM_per_1ms_1e-3\endcsname{yes}
\expandafter\def\csname res@rho_LSTM_per_1ms_3e-3\endcsname{0.837}
\expandafter\def\csname res@rho_LSTM_per_1ms_3e-3_ci\endcsname{[0.827, 0.848]}
\expandafter\def\csname res@sig_LSTM_per_1ms_3e-3\endcsname{yes}
\expandafter\def\csname res@rho_LSTM_per_1ms_1e-2\endcsname{0.542}
\expandafter\def\csname res@rho_LSTM_per_1ms_1e-2_ci\endcsname{[0.527, 0.559]}
\expandafter\def\csname res@sig_LSTM_per_1ms_1e-2\endcsname{yes}
\expandafter\def\csname res@rho_LSTM_per_1ms_3e-2\endcsname{0.248}
\expandafter\def\csname res@rho_LSTM_per_1ms_3e-2_ci\endcsname{[0.232, 0.266]}
\expandafter\def\csname res@sig_LSTM_per_1ms_3e-2\endcsname{yes}
\expandafter\def\csname res@rho_LSTM_per_1ms_1e-1\endcsname{0.080}
\expandafter\def\csname res@rho_LSTM_per_1ms_1e-1_ci\endcsname{[0.071, 0.091]}
\expandafter\def\csname res@sig_LSTM_per_1ms_1e-1\endcsname{yes}
\expandafter\def\csname res@rho_Transformer_per_spike_0\endcsname{1.000}
\expandafter\def\csname res@rho_Transformer_per_spike_0_ci\endcsname{[1.000, 1.000]}
\expandafter\def\csname res@rho_Transformer_per_spike_1e-4\endcsname{1.000}
\expandafter\def\csname res@rho_Transformer_per_spike_1e-4_ci\endcsname{[1.000, 1.000]}
\expandafter\def\csname res@sig_Transformer_per_spike_1e-4\endcsname{no}
\expandafter\def\csname res@rho_Transformer_per_spike_3e-4\endcsname{1.000}
\expandafter\def\csname res@rho_Transformer_per_spike_3e-4_ci\endcsname{[1.000, 1.000]}
\expandafter\def\csname res@sig_Transformer_per_spike_3e-4\endcsname{yes}
\expandafter\def\csname res@rho_Transformer_per_spike_1e-3\endcsname{1.000}
\expandafter\def\csname res@rho_Transformer_per_spike_1e-3_ci\endcsname{[1.000, 1.000]}
\expandafter\def\csname res@sig_Transformer_per_spike_1e-3\endcsname{yes}
\expandafter\def\csname res@rho_Transformer_per_spike_3e-3\endcsname{0.999}
\expandafter\def\csname res@rho_Transformer_per_spike_3e-3_ci\endcsname{[0.999, 1.000]}
\expandafter\def\csname res@sig_Transformer_per_spike_3e-3\endcsname{no}
\expandafter\def\csname res@rho_Transformer_per_spike_1e-2\endcsname{0.998}
\expandafter\def\csname res@rho_Transformer_per_spike_1e-2_ci\endcsname{[0.998, 0.999]}
\expandafter\def\csname res@sig_Transformer_per_spike_1e-2\endcsname{yes}
\expandafter\def\csname res@rho_Transformer_per_spike_3e-2\endcsname{0.994}
\expandafter\def\csname res@rho_Transformer_per_spike_3e-2_ci\endcsname{[0.993, 0.995]}
\expandafter\def\csname res@sig_Transformer_per_spike_3e-2\endcsname{yes}
\expandafter\def\csname res@rho_Transformer_per_spike_1e-1\endcsname{0.980}
\expandafter\def\csname res@rho_Transformer_per_spike_1e-1_ci\endcsname{[0.977, 0.982]}
\expandafter\def\csname res@sig_Transformer_per_spike_1e-1\endcsname{yes}
\expandafter\def\csname res@rho_Transformer_per_1ms_0\endcsname{1.000}
\expandafter\def\csname res@rho_Transformer_per_1ms_0_ci\endcsname{[1.000, 1.000]}
\expandafter\def\csname res@rho_Transformer_per_1ms_1e-4\endcsname{0.995}
\expandafter\def\csname res@rho_Transformer_per_1ms_1e-4_ci\endcsname{[0.994, 0.996]}
\expandafter\def\csname res@sig_Transformer_per_1ms_1e-4\endcsname{yes}
\expandafter\def\csname res@rho_Transformer_per_1ms_3e-4\endcsname{0.985}
\expandafter\def\csname res@rho_Transformer_per_1ms_3e-4_ci\endcsname{[0.984, 0.987]}
\expandafter\def\csname res@sig_Transformer_per_1ms_3e-4\endcsname{yes}
\expandafter\def\csname res@rho_Transformer_per_1ms_1e-3\endcsname{0.949}
\expandafter\def\csname res@rho_Transformer_per_1ms_1e-3_ci\endcsname{[0.945, 0.954]}
\expandafter\def\csname res@sig_Transformer_per_1ms_1e-3\endcsname{yes}
\expandafter\def\csname res@rho_Transformer_per_1ms_3e-3\endcsname{0.836}
\expandafter\def\csname res@rho_Transformer_per_1ms_3e-3_ci\endcsname{[0.826, 0.846]}
\expandafter\def\csname res@sig_Transformer_per_1ms_3e-3\endcsname{yes}
\expandafter\def\csname res@rho_Transformer_per_1ms_1e-2\endcsname{0.547}
\expandafter\def\csname res@rho_Transformer_per_1ms_1e-2_ci\endcsname{[0.530, 0.564]}
\expandafter\def\csname res@sig_Transformer_per_1ms_1e-2\endcsname{yes}
\expandafter\def\csname res@rho_Transformer_per_1ms_3e-2\endcsname{0.281}
\expandafter\def\csname res@rho_Transformer_per_1ms_3e-2_ci\endcsname{[0.266, 0.296]}
\expandafter\def\csname res@sig_Transformer_per_1ms_3e-2\endcsname{yes}
\expandafter\def\csname res@rho_Transformer_per_1ms_1e-1\endcsname{0.114}
\expandafter\def\csname res@rho_Transformer_per_1ms_1e-1_ci\endcsname{[0.098, 0.128]}
\expandafter\def\csname res@sig_Transformer_per_1ms_1e-1\endcsname{yes}
\expandafter\def\csname res@r_bin10ms_s0\endcsname{0.858}
\expandafter\def\csname res@rho_bin10ms_s0_1e-3\endcsname{0.964}
\expandafter\def\csname res@rho_bin10ms_s0_1e-3_ci\endcsname{[0.963, 0.965]}
\expandafter\def\csname res@rho_bin10ms_s0_1e-2\endcsname{0.758}
\expandafter\def\csname res@rho_bin10ms_s0_1e-2_ci\endcsname{[0.750, 0.765]}
\expandafter\def\csname res@r_bin10ms_s3\endcsname{0.869}
\expandafter\def\csname res@rho_bin10ms_s3_1e-3\endcsname{0.970}
\expandafter\def\csname res@rho_bin10ms_s3_1e-3_ci\endcsname{[0.969, 0.972]}
\expandafter\def\csname res@rho_bin10ms_s3_1e-2\endcsname{0.787}
\expandafter\def\csname res@rho_bin10ms_s3_1e-2_ci\endcsname{[0.780, 0.794]}
\expandafter\def\csname res@r_bin20ms_s0\endcsname{0.858}
\expandafter\def\csname res@rho_bin20ms_s0_1e-3\endcsname{0.971}
\expandafter\def\csname res@rho_bin20ms_s0_1e-3_ci\endcsname{[0.969, 0.972]}
\expandafter\def\csname res@rho_bin20ms_s0_1e-2\endcsname{0.788}
\expandafter\def\csname res@rho_bin20ms_s0_1e-2_ci\endcsname{[0.782, 0.795]}
\expandafter\def\csname res@r_bin20ms_s3\endcsname{0.873}
\expandafter\def\csname res@rho_bin20ms_s3_1e-3\endcsname{0.979}
\expandafter\def\csname res@rho_bin20ms_s3_1e-3_ci\endcsname{[0.978, 0.980]}
\expandafter\def\csname res@rho_bin20ms_s3_1e-2\endcsname{0.840}
\expandafter\def\csname res@rho_bin20ms_s3_1e-2_ci\endcsname{[0.833, 0.846]}
\expandafter\def\csname res@r_bin50ms_s0\endcsname{0.853}
\expandafter\def\csname res@rho_bin50ms_s0_1e-3\endcsname{0.960}
\expandafter\def\csname res@rho_bin50ms_s0_1e-3_ci\endcsname{[0.957, 0.962]}
\expandafter\def\csname res@rho_bin50ms_s0_1e-2\endcsname{0.731}
\expandafter\def\csname res@rho_bin50ms_s0_1e-2_ci\endcsname{[0.724, 0.738]}
\expandafter\def\csname res@r_bin50ms_s3\endcsname{0.844}
\expandafter\def\csname res@rho_bin50ms_s3_1e-3\endcsname{0.979}
\expandafter\def\csname res@rho_bin50ms_s3_1e-3_ci\endcsname{[0.977, 0.981]}
\expandafter\def\csname res@rho_bin50ms_s3_1e-2\endcsname{0.843}
\expandafter\def\csname res@rho_bin50ms_s3_1e-2_ci\endcsname{[0.836, 0.850]}
\expandafter\def\csname res@r_N20\endcsname{0.572}
\expandafter\def\csname res@rho_N20_1e-3\endcsname{0.952}
\expandafter\def\csname res@rho_N20_1e-2\endcsname{0.713}
\expandafter\def\csname res@r_N40\endcsname{0.658}
\expandafter\def\csname res@rho_N40_1e-3\endcsname{0.936}
\expandafter\def\csname res@rho_N40_1e-2\endcsname{0.651}
\expandafter\def\csname res@r_N80\endcsname{0.792}
\expandafter\def\csname res@rho_N80_1e-3\endcsname{0.957}
\expandafter\def\csname res@rho_N80_1e-2\endcsname{0.751}
\expandafter\def\csname res@r_N161\endcsname{0.858}
\expandafter\def\csname res@rho_N161_1e-3\endcsname{0.965}
\expandafter\def\csname res@rho_N161_1e-2\endcsname{0.748}
\expandafter\def\csname res@pr\endcsname{14.1}
\expandafter\def\csname res@pr_shuf\endcsname{20.8}
}{}
\IfFileExists{results/derived.tex}{% derived from results.json (absolute r_SER = mean over folds)
\expandafter\def\csname res@rabs_Ridge_1ms_1e-3\endcsname{0.828}
\expandafter\def\csname res@rabs_Ridge_1ms_3e-3\endcsname{0.774}
\expandafter\def\csname res@rabs_Ridge_1ms_1e-2\endcsname{0.649}
\expandafter\def\csname res@rabs_MLP_1ms_1e-3\endcsname{0.855}
\expandafter\def\csname res@rabs_MLP_1ms_3e-3\endcsname{0.760}
\expandafter\def\csname res@rabs_MLP_1ms_1e-2\endcsname{0.546}
\expandafter\def\csname res@rabs_LSTM_1ms_1e-3\endcsname{0.855}
\expandafter\def\csname res@rabs_LSTM_1ms_3e-3\endcsname{0.756}
\expandafter\def\csname res@rabs_LSTM_1ms_1e-2\endcsname{0.489}
\expandafter\def\csname res@rabs_Transformer_1ms_1e-3\endcsname{0.851}
\expandafter\def\csname res@rabs_Transformer_1ms_3e-3\endcsname{0.750}
\expandafter\def\csname res@rabs_Transformer_1ms_1e-2\endcsname{0.490}
}{}
\newcommand{\todo}[1]{\textcolor{orange!70!black}{\textit{[To write after the rerun: #1]}}}
\newcommand{\fig}[2]{\IfFileExists{results/#1}{\includegraphics[width=#2\linewidth]{results/#1}}%
  {\fbox{\parbox[c][4cm][c]{#2\linewidth}{\centering\small\texttt{results/\detokenize{#1}}\par(run the notebook)}}}}

\title{\vspace{-1.2cm}\textbf{How much spike-transmission error can a rate-based motor decoder tolerate?}\\[3pt]
\large Much less than a per-spike error model suggests, and it is partly set by temporal integration}
\author{Jihun Yu \quad {\small Dept.\ of Bio and Brain Engineering, KAIST \quad \texttt{youjhun0225@kaist.ac.kr}}}
\date{\small Draft v3, October 2026 \quad Code: \url{https://github.com/youjhun/spike-error-tolerance-decoders}}

\begin{document}
\maketitle
\vspace{-0.6cm}

% ════════════════════════════════ ONE-PAGE SUMMARY ════════════════════════════════
\section*{Summary}

\textbf{Question.} Lee et al.\ (2024) showed that an ASBIT wireless link keeps 88.8\% (M1) and 85.4\% (S1) of
wired decoding accuracy at a spike error rate (SER) of $10^{-3}$, using an SNN decoder on 1-ms data, and
suggested that real populations, which share information, should be even more robust~\cite{lee2024}.
I asked how much SER classical rate-based decoders tolerate on a public macaque dataset, and why.

\textbf{Setup.} CRCNS pmd-1 (\res{n_m1} M1 + \res{n_pmd} PMd units, \res{n_trials} trials, 10-ms bins)~\cite{pmd1};
Ridge, MLP, LSTM and causal Transformer decoders of 2-D hand velocity from a 200-ms causal window;
trial-level 5-fold cross-validation; decoders trained once on clean data and tested on corrupted data.
Spike errors were injected under two definitions of SER (Sec.~\ref{sec:ser}).

\begin{figure}[H]\centering \fig{fig1_ser_definitions.png}{0.82}
\caption{\textbf{Left:} retention $\rho=r_{\mathrm{SER}}/r_{\mathrm{clean}}$ (mean, 95\% bootstrap CI) for each decoder.
Dashed: SER normalised per spike (used in my previous report). Solid: SER as the fraction of 1-ms bins in error
(my reading of Lee et al.). Crosses: Lee et al.'s SNN at SER $10^{-3}$ (different data and decoder; shown for scale only).
\textbf{Right:} false spikes each definition adds per neuron, against the mean true rate (\res{rate_hz}\,Hz).}
\end{figure}

\textbf{Results.}
\begin{enumerate}
  \item \textbf{Two SER definitions.} At SER $=10^{-2}$, Ridge retains $\rho=$ \res{rho_Ridge_per_spike_1e-2}
  when errors are counted per spike, but \res{rho_Ridge_per_1ms_1e-2} \res{rho_Ridge_per_1ms_1e-2_ci} when they are counted per
  1-ms bin. The per-bin definition adds \res{fhz_1ms_1e-2}\,Hz of false spikes per neuron, versus \res{fhz_spike_1e-2}\,Hz
  per spike. Under the per-bin model the more flexible decoders lose more: $\rho=$ \res{rho_MLP_per_1ms_1e-2} (MLP), \res{rho_LSTM_per_1ms_1e-2} (LSTM) and \res{rho_Transformer_per_1ms_1e-2} (Transformer). They are the most accurate on clean data, but from SER $3\times10^{-3}$ on Ridge gives the highest absolute $r$ (\res{rabs_Ridge_1ms_3e-3} vs \res{rabs_MLP_1ms_3e-3} for MLP).
  \item \textbf{Temporal integration.} Summing to 20-ms bins and smoothing the output raised Ridge retention at SER $10^{-2}$ from \res{rho_bin10ms_s0_1e-2} to \res{rho_bin20ms_s3_1e-2} (baseline re-estimated with new corruption draws), with no loss of clean accuracy ($r=$ \res{r_bin10ms_s0} $\to$ \res{r_bin20ms_s3}; Fig.~\ref{fig:int}).
  \item \textbf{Population size.} Retention did not rise with the number of units (SER $10^{-2}$: \res{rho_N20_1e-2} at $N=20$, \res{rho_N161_1e-2} at $N=161$; Fig.~\ref{fig:pop}), although clean accuracy did. The ``population dilution'' explanation of my previous report is not supported.
\end{enumerate}

\textbf{Reading against Lee et al.} Under the per-bin model, the rate-based decoders here keep 94--96\% at SER $10^{-3}$, against 88.8\% for their SNN. Data, channel count and decoder differ, so this gap is not a measured advantage.

\textbf{What this changes.} My previous report claimed a 30-point robustness advantage of rate-based decoders over the
SNN. That comparison mixed two SER definitions and misdescribed the SNN, which itself bins spikes into 25-ms
intervals~\cite{lee2024}; it is withdrawn (Appendix~\ref{app:changes}).

\textbf{Limits.} One session of one monkey; i.i.d.\ errors; no SNN or ASBIT channel simulation; 10-ms data cannot
reproduce 1-ms effects exactly.

\newpage
% ════════════════════════════════ MAIN TEXT ════════════════════════════════
\section{Motivation}
In the ASBIT protocol each sensor sends a Gold-coded packet only when it detects an event, so link errors appear as
missed or false spikes rather than corrupted samples~\cite{lee2024}. The event error rate (EER) a link must reach
therefore depends on how many spike errors the downstream decoder can absorb. Lee et al.\ found that network
performance is set mainly by the total number of events per second (about $10^5$ events/s at EER $<10^{-3}$), and
that the number of nodes is likely limited by wireless power transfer and SAR rather than bandwidth~\cite{lee2024}.
A decoder that tolerates more spike errors would relax the EER target. This report measures that tolerance for
rate-based decoders and asks which property of the decoder explains it.

\section{Methods}
\paragraph{Data.} CRCNS pmd-1, file \texttt{MM\_S1\_processed.mat} (``S1'' denotes session 1,
not somatosensory cortex)~\cite{pmd1}. \res{n_neurons} units (\res{n_m1} M1, \res{n_pmd} PMd), \res{n_trials} trials,
10-ms bins; \res{sparsity}\% of unit-bins are empty, mean rate \res{rate_hz}\,Hz, \res{kbar} active units per bin.

\paragraph{Features and validation.} For each trial separately, $\mathbf{x}_t=[\mathbf{s}_{t-19},\dots,\mathbf{s}_t]$
(20 bins $\times$ \res{n_neurons} units). Folds are split by trial; lag windows never cross trial boundaries.
No spike-count transform and no output smoothing are used in the main analysis; smoothing ($\sigma=3$ bins,
applied within each trial) appears only in Sec.~\ref{sec:int}.

\paragraph{Decoders.} Ridge (inputs standardised; penalty chosen by inner 3-fold trial-level CV from
$\{1,10,\dots,10^4\}$); MLP ($3220\!\to\!512\!\to\!256\!\to\!2$, dropout 0.2); 2-layer LSTM (hidden 128);
3-layer causal Transformer ($d=128$, 4 heads). Neural networks use Adam ($10^{-3}$, weight decay $10^{-4}$) with
early stopping on 10\% of training trials. Accuracy is Pearson $r$ averaged over $v_x$ and $v_y$.

\paragraph{Spike-error models.}\label{sec:ser}
Decoders are trained once on clean data; only test trials are corrupted, with \texttt{N\_CORRUPT\_REPEATS}~=~5
independent draws per fold and SER level.
\begin{itemize}
  \item \textbf{Per spike} (previous report): $\mathrm{SER}=(N_{\mathrm{miss}}+N_{\mathrm{false}})/N_{\mathrm{spike}}$;
  each spike is deleted with probability $p/2$ and $\mathrm{Binomial}(N_{\mathrm{spike}},p/2)$ false spikes are placed at random.
  \item \textbf{Per 1-ms bin} (my reading of Lee et al., who count a missed or false spike within a 1-ms bin as an
  error~\cite{lee2024}; the denominator is not stated explicitly): each 10-ms bin is treated as ten 1-ms slots; a
  spike is missed with probability $p$, and each empty slot receives a false spike with probability $p$.
\end{itemize}
Under the per-bin model a neuron firing at $\lambda$\,Hz gains about $1000p$ false spikes per second; under the
per-spike model it gains $\lambda p/2$. At the rates in this dataset the two differ by more than two orders of magnitude.

\paragraph{Statistics.} $\rho=r_{\mathrm{SER}}/r_{\mathrm{clean}}$ per fold and draw; 95\% bootstrap CIs over these
values (not independent across draws within a fold; see Limitations). A drop is tested by a paired $t$-test on Fisher-$z$ transformed per-fold correlations (the test Lee et al.\ used for their jitter comparison),
with Benjamini--Hochberg correction across SER levels.

\section{Results}
\subsection{Clean accuracy}
\begin{table}[H]\centering\small
\begin{tabular}{lcccc}\toprule
 & Ridge & MLP & LSTM & Transformer\\\midrule
$r$ (mean $v_x$, $v_y$) & \res{r_Ridge} & \res{r_MLP} & \res{r_LSTM} & \res{r_Transformer}\\
95\% CI & \res{r_Ridge_ci} & \res{r_MLP_ci} & \res{r_LSTM_ci} & \res{r_Transformer_ci}\\
$r$ ($v_x$ only) & \res{rvx_Ridge} & \res{rvx_MLP} & \res{rvx_LSTM} & \res{rvx_Transformer}\\
$r$ with smoothing $\sigma=3$ & \res{rsm_Ridge} & \res{rsm_MLP} & \res{rsm_LSTM} & \res{rsm_Transformer}\\\bottomrule
\end{tabular}
\caption{Clean decoding, trial-level 5-fold CV. Lee et al.\ report $r=0.9324$ for $v_x$ with an SNN on different M1
data; the numbers are not directly comparable.}\end{table}

\subsection{Retention under the two SER definitions}
\begin{table}[H]\centering\small
\begin{tabular}{lcccccc}\toprule
 & \multicolumn{2}{c}{SER $10^{-3}$} & \multicolumn{2}{c}{SER $10^{-2}$} & \multicolumn{2}{c}{SER $10^{-1}$}\\
\cmidrule(lr){2-3}\cmidrule(lr){4-5}\cmidrule(lr){6-7}
Decoder & per spike & per 1 ms & per spike & per 1 ms & per spike & per 1 ms\\\midrule
Ridge & \res{rho_Ridge_per_spike_1e-3} & \res{rho_Ridge_per_1ms_1e-3} & \res{rho_Ridge_per_spike_1e-2} & \res{rho_Ridge_per_1ms_1e-2} & \res{rho_Ridge_per_spike_1e-1} & \res{rho_Ridge_per_1ms_1e-1}\\
MLP & \res{rho_MLP_per_spike_1e-3} & \res{rho_MLP_per_1ms_1e-3} & \res{rho_MLP_per_spike_1e-2} & \res{rho_MLP_per_1ms_1e-2} & \res{rho_MLP_per_spike_1e-1} & \res{rho_MLP_per_1ms_1e-1}\\
LSTM & \res{rho_LSTM_per_spike_1e-3} & \res{rho_LSTM_per_1ms_1e-3} & \res{rho_LSTM_per_spike_1e-2} & \res{rho_LSTM_per_1ms_1e-2} & \res{rho_LSTM_per_spike_1e-1} & \res{rho_LSTM_per_1ms_1e-1}\\
Transformer & \res{rho_Transformer_per_spike_1e-3} & \res{rho_Transformer_per_1ms_1e-3} & \res{rho_Transformer_per_spike_1e-2} & \res{rho_Transformer_per_1ms_1e-2} & \res{rho_Transformer_per_spike_1e-1} & \res{rho_Transformer_per_1ms_1e-1}\\\midrule
False spikes added (Hz) & \res{fhz_spike_1e-3} & \res{fhz_1ms_1e-3} & \res{fhz_spike_1e-2} & \res{fhz_1ms_1e-2} & \res{fhz_spike_1e-1} & \res{fhz_1ms_1e-1}\\\bottomrule
\end{tabular}
\caption{Retention $\rho$ (mean over 5 folds $\times$ 5 draws). Full sweep, CIs and BH-corrected tests are in
\texttt{results.json}.}\end{table}
Under the per-spike model, retention stays above 0.97 for every decoder up to SER $10^{-1}$. Under the per-1-ms model it falls to \res{rho_MLP_per_1ms_1e-3}--\res{rho_Ridge_per_1ms_1e-3} at $10^{-3}$ and to \res{rho_LSTM_per_1ms_1e-2}--\res{rho_Ridge_per_1ms_1e-2} at $10^{-2}$; Ridge degrades least and LSTM most. Every per-1-ms drop from $10^{-4}$ upward survives Benjamini--Hochberg correction; per-spike drops become statistically detectable between $10^{-3}$ and $10^{-2}$, depending on the decoder, but stay at or below 0.3\% up to $10^{-2}$. Ridge's penalty was chosen at the top of the grid ($\alpha=10^4$, recorded for the last fold), so heavier shrinkage may raise its tolerance further.

\subsection{Temporal integration}\label{sec:int}
If binning acts as a low-pass filter on spike trains, wider bins should average out uncorrelated spike errors,
whose power is spread across frequencies, while keeping the slow velocity-related signal. I re-binned the
(already corrupted) 10-ms counts to 20 and 50\,ms, keeping a 200-ms window, with and without output smoothing.
\begin{figure}[H]\centering \fig{fig2_integration.png}{0.62}
\caption{Ridge retention under per-1-ms errors for different bin widths and smoothing ($\sigma$ in bins).}\label{fig:int}\end{figure}
Going from 10- to 20-ms bins and adding output smoothing each raised retention at SER $10^{-2}$; together, 20-ms bins with $\sigma=3$ gave \res{rho_bin20ms_s3_1e-2} \res{rho_bin20ms_s3_1e-2_ci}, against \res{rho_bin10ms_s0_1e-2} for the 10-ms baseline. Clean accuracy was unchanged or higher (\res{r_bin10ms_s0}, \res{r_bin20ms_s3}). At 50 ms without smoothing retention fell to \res{rho_bin50ms_s0_1e-2}, likely because only 4 lags remain; with smoothing it was \res{rho_bin50ms_s3_1e-2} but clean $r$ dropped to \res{r_bin50ms_s3}. Low-pass filtering helps, but recovers only part of the loss.

\subsection{Population size and redundancy}
\begin{figure}[H]\centering \fig{fig3_population.png}{0.55}
\caption{Ridge retention under per-1-ms errors for random subsets of $N$ units (5 subsets each).}\label{fig:pop}\end{figure}
The participation ratio of smoothed population activity is \res{pr} (out of \res{n_neurons} units), against \res{pr_shuf}
after shuffling time bins within trials. So population activity is low-dimensional relative to its size, and more so than after destroying temporal structure. Yet adding units did not improve retention (Fig.~\ref{fig:pop}; CIs overlap from $N=20$ to 161). One reading is that each added unit brings its own independent false spikes as well as signal, so redundancy improves clean accuracy without improving the ratio. Testing this needs a population-structure manipulation, not random subsampling.

\section{Discussion}
If the per-1-ms model is the right reading, rate-based decoders on this dataset need SER of about $10^{-3}$ or lower to keep roughly 95\% of clean accuracy (Ridge \res{rho_Ridge_per_1ms_1e-3}; neural networks \res{rho_MLP_per_1ms_1e-3}--\res{rho_Transformer_per_1ms_1e-3}). This is close to the EER of about $10^{-3}$ that Lee et al.\ describe as acceptable for BMI use. Decoder-side measures (wider bins, smoothing, heavier shrinkage) buy only part of a decade. This suggests, as a hypothesis, that most of the error budget will have to come from the link and its power delivery, rather than from choosing a more tolerant decoder.

A decoder's tolerance only matters relative to the error model of a real link. The per-bin model here assumes
independent errors at a constant rate; Lee et al.\ show that a node's error rate depends mostly on its own
backscatter amplitude, so near--far differences in SNR would make errors unequal across channels~\cite{lee2024}.
Testing decoders under that channel-dependent error pattern is the natural next step, and it connects decoder
tolerance to the power and SAR budget that limits node count.

\section{Limitations}
\begin{itemize}
  \item One monkey, one session; no cross-session test.
  \item The data are 10-ms bins, so 1-ms errors are emulated by sub-slot sampling.
  \item No SNN was trained and no ASBIT link was simulated; comparison with Lee et al.\ is qualitative only.
  \item The SER denominator in Lee et al.\ is my reading of the text and should be confirmed.
  \item Ridge's penalty hit the top of the search grid; a wider grid is needed.
  \item Confidence intervals pool fold $\times$ draw values (25 per condition); draws within a fold share one trained decoder and one test set, so they are not independent and the intervals are optimistic.
\end{itemize}

\appendix
\section{Changes from the previous version (v2, June 2026)}\label{app:changes}
\begin{itemize}
  \item \textbf{Withdrawn: ``30-point advantage over the SNN''.} v2 compared retention under a per-spike SER with the
  SNN's retention, whose SER is (as I read it) counted per 1-ms bin, on different data. The per-spike model adds far fewer false spikes.
  \item \textbf{Corrected: the SNN does not use precise inter-spike intervals.} It bins 1-ms spikes into 25-ms intervals
  and is trained with $\pm$25\,ms jitter~\cite{lee2024}; v2's timing-versus-rate explanation was wrong.
  \item \textbf{Corrected attribution.} The ASBIT paper is Lee, Lee et al.\ (2024), not ``Nurmikko et al.'';
  ASBIT is the network protocol, not the decoder.
  \item \textbf{Removed unsupported claims}: typical link BER values and power savings attributed to the paper.
  \item \textbf{Analysis}: dropped the spike-count exponent $\alpha$ (chosen on a test fold in v2); removed default
  output smoothing; corruption repeated 5 times instead of once; added CIs and BH correction.
  \item \textbf{Dataset}: pmd-1 is from Perich, Lawlor, Kording and Miller and involves sequential reaching~\cite{pmd1};
  v2 cited it incorrectly.
  \item \textbf{Attention analysis} moved out of the main text; the attention weights in v2 (about 0.06--0.07 per bin
  against a uniform 0.05) do not support a strong recency claim.
\end{itemize}

\begin{thebibliography}{9}\small
\bibitem{lee2024} J.~Lee, A.-H.~Lee, V.~Leung, F.~Laiwalla, M.~A.~Lopez-Gordo, L.~Larson and A.~Nurmikko,
``An asynchronous wireless network for capturing event-driven data from large populations of autonomous sensors,''
\textit{Nature Electronics} \textbf{7}, 313--324 (2024). \href{https://doi.org/10.1038/s41928-024-01134-y}{doi:10.1038/s41928-024-01134-y}
\bibitem{pmd1} M.~G.~Perich, P.~N.~Lawlor, K.~P.~Kording and L.~E.~Miller, ``Extracellular neural recordings from
macaque primary and dorsal premotor motor cortex during a sequential reaching task,'' CRCNS.org (pmd-1).
\href{https://doi.org/10.6080/K0FT8J72}{doi:10.6080/K0FT8J72}
\end{thebibliography}
\end{document}
